
\documentclass{article}
\usepackage{colortbl}
\usepackage{makecell}
\usepackage{multirow}
\usepackage{supertabular}

\begin{document}

\newcounter{utterance}

\centering \large Interaction Transcript for game `cladder', experiment `full\_v1.5\_default', episode 4237 with Qwen3.5{-}9B.
\vspace{24pt}

{ \footnotesize  \setcounter{utterance}{1}
\setlength{\tabcolsep}{0pt}
\begin{supertabular}{c@{$\;$}|p{.15\linewidth}@{}p{.15\linewidth}p{.15\linewidth}p{.15\linewidth}p{.15\linewidth}p{.15\linewidth}}
    \# & \multicolumn{2}{c}{Player} && \multicolumn{2}{c}{Game Master} \\
    \hline

    \theutterance \stepcounter{utterance}  
    & & & \multicolumn{4}{p{0.6\linewidth}}{
        \cellcolor[rgb]{0.9,0.9,0.9}{
            \makecell[{{p{\linewidth}}}]{
                \texttt{\tiny{[P1$\langle$GM]}}
                \texttt{You are an expert in causal inference. The following question is not a typical commonsense query, but rather a meticulously designed question created by a professor specializing in causal inference, intended to assess the students' mastery of the course content.} \\
\texttt{Imagine a self{-}contained, hypothetical world with only the following conditions, and without any unmentioned factors or causal relationships: Kidney stone size has a direct effect on treatment and lip thickness. Treatment has a direct effect on lip thickness. For patients who have small kidney stones and not receiving treatment, the probability of thick lips is 6\%. For patients who have small kidney stones and receiving treatment, the probability of thick lips is 38\%. For patients who have large kidney stones and not receiving treatment, the probability of thick lips is 63\%. For patients who have large kidney stones and receiving treatment, the probability of thick lips is 95\%. The overall probability of large kidney stone is 50\%. For patients receiving treatment, would it be more likely to see thick lips if the patient had received no treatment?} \\
\texttt{Start your answer with 'Yes' or 'No', followed by additional reasoning or evidence to support your explanation.} \\
            }
        }
    }
     \\ \\

    \theutterance \stepcounter{utterance}  
    & \multicolumn{4}{p{0.6\linewidth}}{
        \cellcolor[rgb]{0.561,0.737,0.561}{
            \makecell[{{p{\linewidth}}}]{
                \texttt{\tiny{[P1$\rangle$GM]}}
                \texttt{No} \\
\\ 
\texttt{**Reasoning:**} \\
\\ 
\texttt{To answer this question, we must determine the causal effect of treatment on lip thickness. The question asks: "For patients receiving treatment, would it be more likely to see thick lips if the patient had received no treatment?" This is asking for the counterfactual comparison: Is \$P(\textbackslash{}text\{Thick Lips\} | \textbackslash{}text\{Treatment\}=1) > P(\textbackslash{}text\{Thick Lips\} | \textbackslash{}text\{Treatment\}=0, \textbackslash{}text\{Kidney Stone Size\})\$?} \\
\\ 
\texttt{However, the provided probabilities are conditional on both kidney stone size and treatment status (observed data), not on the counterfactual treatment status. To make a valid causal comparison, we need to know if the effect of treatment is consistent across different kidney stone sizes (i.e., whether the **Consistency** and **Homogeneity** assumptions hold, or if we can calculate the Average Treatment Effect).} \\
\\ 
\texttt{Let's analyze the data provided:} \\
\texttt{1.  **Small Kidney Stones:**} \\
\texttt{    *   \$P(\textbackslash{}text\{Thick\} | \textbackslash{}text\{Small\}, \textbackslash{}text\{Treat\}=0) = 0.06\$} \\
\texttt{    *   \$P(\textbackslash{}text\{Thick\} | \textbackslash{}text\{Small\}, \textbackslash{}text\{Treat\}=1) = 0.38\$} \\
\texttt{    *   Effect of Treatment (Small) = \$0.38 {-} 0.06 = +0.32\$ (Treatment} \\
            }
        }
    }
    & & \\ \\

    \theutterance \stepcounter{utterance}  
    & & & \multicolumn{2}{p{0.3\linewidth}}{
        \cellcolor[rgb]{0.9,0.9,0.9}{
            \makecell[{{p{\linewidth}}}]{
                \texttt{\tiny{[GM$|$GM]}}
                \texttt{No} \\
            }
        }
    }
    & & \\ \\

    \theutterance \stepcounter{utterance}  
    & & & \multicolumn{2}{p{0.3\linewidth}}{
        \cellcolor[rgb]{0.9,0.9,0.9}{
            \makecell[{{p{\linewidth}}}]{
                \texttt{\tiny{[GM$|$GM]}}
                \texttt{no} \\
            }
        }
    }
    & & \\ \\

    \theutterance \stepcounter{utterance}  
    & & & \multicolumn{2}{p{0.3\linewidth}}{
        \cellcolor[rgb]{0.9,0.9,0.9}{
            \makecell[{{p{\linewidth}}}]{
                \texttt{\tiny{[GM$|$GM]}}
                \texttt{game\_result = WIN} \\
            }
        }
    }
    & & \\ \\

\end{supertabular}
}

\end{document}
